\chapter{Conclusions and Future Work}
\label{chap6}

\section{Conclusions for the Research Questions}
\label{sec:conclusion-research-questions}

This dissertation asked whether conversational context and speech provide useful affective evidence, how such evidence can be used without giving a classifier authority over a role-play, and whether the resulting application meets its declared engineering requirements. The evaluation in Chapter~\ref{chap5} supports the following bounded answers.

\paragraph{RQ1: conversational context.}
\label{subsec:conclusion-rq1}
Preceding dialogue turns improved text-based emotion classification under the declared four-class, speaker-independent IEMOCAP protocol. The paired dialogue-level comparison consistently favoured the contextual model, so the improvement cannot be attributed to different speakers, label sets, or optimisation settings. This is evidence that causal conversational context helps classify annotated expression in this benchmark; it is not evidence that the model recovers a person's private emotional state or that the gain will transfer unchanged to spontaneous AffectLab conversations. The exploratory six-class result points in the same aggregate direction but is too sparse to support an anxiety-recognition claim.

\paragraph{RQ2: multimodal complementarity and calibration.}
\label{subsec:conclusion-rq2}
Acoustic information was complementary to contextual text even though audio was the weaker individual modality. Predeclared equal-weight fusion improved the principal classification measures over either unimodal model, with the paired bootstrap comparison favouring fusion. However, this configuration was substantially less calibrated than contextual text. Validation-fitted fusion recovered strong calibration while retaining most of the classification gain. The answer is therefore conditional: multimodal fusion improved benchmark classification, but the configuration that maximised classification was not the configuration with the most trustworthy probabilities. Classification and calibration must be selected and reported separately when confidence affects adaptive behaviour.

\paragraph{RQ3: bounded use of affective evidence.}
\label{subsec:conclusion-rq3}
Uncertain predictions can be integrated responsibly by restricting what they are allowed to change. AffectLab records the prediction and its provenance, checks confidence, modality agreement, availability, eligibility, and user preference, and maps the result to a small set of reversible pacing actions. Deterministic scenario controllers remain authoritative for progression and outcomes; observable language evidence determines feedback; explicit safety handling takes precedence; and inference failure preserves the baseline workflow. Authored policy cases and integrated journeys conformed to these rules. This establishes explainable and failure-tolerant policy behaviour under tested conditions, not that an affect estimate is correct for a particular user or that adaptation improves rehearsal outcomes.

\paragraph{RQ4: integrated application quality.}
\label{subsec:conclusion-rq4}
The completed application satisfied its scoped functional and non-functional requirements for the tested offline deterministic core. Backend, production-frontend, browser, automated-accessibility, persistence, recovery, and full-stack checks supported the declared behaviours, while the local regression benchmark met every response-time budget. These layers establish an implementation-complete, deployment-oriented system rather than deployment readiness. The performance run excluded real network transit, production persistence, trained inference, external providers, concurrency, and adversarial traffic. A controlled deployment therefore remains conditional on target-environment performance, security, observability, backup recovery, and incident-response validation.

\section{Contributions and Material Limitations}
\label{sec:conclusion-contributions}

The work makes three connected contributions. Empirically, it provides a controlled speaker-independent comparison of utterance-only and contextual text classification, followed by a text--audio study that treats calibration as distinct from classification. Architecturally, it demonstrates a hybrid role-play system in which probabilistic models and optional generation influence presentation without controlling dialogue outcomes, assessment, or safety. Methodologically, it separates benchmark evidence, deterministic policy conformance, application verification, local performance, and supplementary user feedback so that one evidence type is not used to support a stronger claim than it can bear. Versioned policies, provenance, evidence-grounded feedback, fallbacks, and replayable state make the implemented decisions inspectable.

These contributions are bounded by the available data and evaluation. IEMOCAP contains ten actors in acted or improvised English dialogue and does not provide external, cultural, device, or deployment validation. A single primary training seed leaves optimisation variability incompletely characterised, while full-data packages and the operational fusion summary have no independent held-out estimate. Emotion categories and confidence thresholds are operational signals rather than psychological measurements. The three scenario policies and lexical feedback rules cover limited forms of communication, and the separately implemented cognitive support-routing fields are unvalidated heuristics, not measures of cognition, intention, readiness, resistance, or clinical state.

The engineering tests used controlled inputs and synthetic accounts, and automated accessibility checks do not replace evaluation with assistive-technology users. Most importantly, the small, uncontrolled pilot exercised the deterministic baseline rather than multimodal inference or affect-aware pacing. It can inform workflow refinement but cannot establish adaptation benefit, sustained learning, transfer, or absence of harm. AffectLab is consequently a communication-practice research application, not a diagnostic, clinical, crisis, or therapeutic system.

\section{Prioritised Future Work}
\label{sec:conclusion-future-work}

The immediate priority is deployment-specific validation of the existing system. A staging study should exercise the packaged models, external providers, intended database, real network path, and representative concurrency. Declared service objectives, load and soak tests, provider and database failures, backup restoration, security review, monitoring, and incident exercises should establish latency, capacity, recovery, and operational risk before any production-readiness claim.

The next priority is to strengthen the evidence behind both inference and interaction. Cross-corpus and target-context evaluation should examine calibration under language, microphone, noise, transcription, and demographic variation, using repeated seeds and prospective held-out data. Independent human annotation should assess the validity and error patterns of feedback rules; psychologically suggestive support-routing labels should be removed, renamed as observable behaviour, or subjected to appropriate expert and ethics-led validation. A preregistered comparison of the baseline and bounded affect-aware policy should then measure helpfulness, intrusiveness, task outcomes, and user preference without treating engagement as benefit.

Longer-term work should use participatory accessibility testing, multilingual and cultural adaptation, and longitudinal studies of retention and transfer. These studies should preserve opt-in controls, data minimisation, transparent uncertainty, and a useful non-adaptive baseline. Model, prompt, scenario, threshold, and safety-resource changes should remain versioned and governed as the evidence base develops.

\section{Final Conclusion}
\label{sec:final-conclusion}

This dissertation shows that contextual and acoustic affective evidence can improve benchmark classification and can be connected to role-play without making emotion recognition authoritative. AffectLab combines bounded adaptation with deterministic progression, evidence-grounded feedback, safety precedence, user control, and auditable fallbacks. Its evaluation establishes the implemented system and policies within clearly defined test conditions, while leaving claims about emotional truth, human benefit, and production readiness to future evidence. Within those limits, the work provides a reproducible foundation for studying emotion-aware communication rehearsal while preserving uncertainty and user agency.
